\documentclass[10pt,a4paper]{article}
\usepackage[utf8]{inputenc}
\usepackage[spanish,es-noquot]{babel}
\usepackage{amsmath, amssymb}
\usepackage[table,xcdraw]{xcolor}
\usepackage{tabularx}
\usepackage{array}

% Definición de colores parecidos a los de la imagen
\definecolor{headergreen}{HTML}{B8D89C}
\definecolor{pinkcell}{HTML}{F4DCDC}
\definecolor{bluecell}{HTML}{E2F0D9} % o ligero tono azul según la fila

\setlength{\arrayrulewidth}{1pt}

\begin{document}

\begin{center}
{\huge \textbf{\textcolor{green!60!black}{\underline{RESUMEN}}}}
\vspace{0.5cm}

\renewcommand{\arraystretch}{1.5}


\begin{tabularx}{\textwidth}{|X|X|}
\hline
Producto y cociente de potencias & En el producto de potencias con la misma base se suman los exponentes. \\ \hline
Potencia de un producto y de un cociente & La potencia de un producto es igual al producto de las potencias de cada uno de los factores. \\ \hline
\end{tabularx}

\begin{tabular}{|c|c|c|}
    \hline

    \rowcolor{pinkcell!40}
    Producto y cociente de potencias & 
    En el producto de potencias con la misma base se suman los exponentes. En el cociente se restan los exponentes. \newline
    Con el mismo exponente: En el producto, se multiplican las bases y se eleva el resultado al mismo exponente. En el cociente se dividen las bases y se eleva el resultado al mismo exponente. & 
    $(-5)^4 \cdot (-5)^2 = (-5)^6$ \newline
    $3^2 : 3^7 = 3^{-5}$ \newline
    $2^5 \cdot 7^5 = 14^5$ \newline
    $(-5)^3 : (4)^3 = (-5/4)^3$ \\ \hline

    \rowcolor{pinkcell!40}
    Potencia de un producto y de un cociente & 
    La potencia de un producto es igual al producto de cada uno de los factores elevados a dicha potencia:
    $ (a \cdot b \cdot c \cdot d)^n = a^n \cdot b^n \cdot c^n \cdot d^n $
    La potencia de un cociente es igual al cociente del dividendo y el divisor elevados a dicha potencia:
    $ c^m : c^n = c^{m-n} $ & 
    $(5 \cdot 2 \cdot 3)^4 = 5^4 \cdot 2^4 \cdot 3^4$ \vspace{0.3cm}\newline
    $(-7/2)^6 = 7^6 / (-2)^6$ \\ \hline





\textbf{Columna 1} & \textbf{Columna 2} & \textbf{Columna 3} \\ \hline
Dato 1             & Dato 2             & Dato 3             \\ 
Dato 4             & Dato 5             & Dato 6             \\ \hline
\end{tabular}


\begin{tabularx}{\textwidth}{|>{\raggedright\arraybackslash\bfseries}p{3.5cm}|X|>{\centering\arraybackslash}p{4cm}|}
\hline
\rowcolor{headergreen}
 & \centering \textbf{POTENCIAS Y RAÍCES} & \textit{\textbf{Ejemplos}} \\ \hline

\rowcolor{pinkcell!40}
Producto y cociente de potencias & 
En el producto de potencias con la misma base se suman los exponentes. En el cociente se restan los exponentes. \newline
Con el mismo exponente: En el producto, se multiplican las bases y se eleva el resultado al mismo exponente. En el cociente se dividen las bases y se eleva el resultado al mismo exponente. & 
$(-5)^4 \cdot (-5)^2 = (-5)^6$ \newline
$3^2 : 3^7 = 3^{-5}$ \newline
$2^5 \cdot 7^5 = 14^5$ \newline
$(-5)^3 : (4)^3 = (-5/4)^3$ \\ \hline

\rowcolor{pinkcell!40}
Potencia de un producto y de un cociente & 
La potencia de un producto es igual al producto de cada uno de los factores elevados a dicha potencia:
\[ (a \cdot b \cdot c \cdot d)^n = a^n \cdot b^n \cdot c^n \cdot d^n \]
La potencia de un cociente es igual al cociente del dividendo y el divisor elevados a dicha potencia:
\[ c^m : c^n = c^{m-n} \] & 
$(5 \cdot 2 \cdot 3)^4 = 5^4 \cdot 2^4 \cdot 3^4$ \vspace{0.3cm}\newline
$(-7/2)^6 = 7^6 / (-2)^6$ \\ \hline

\rowcolor{pinkcell!40}
Potencia de otra potencia & 
\centering $((d)^m)^n = d^{m \cdot n}$ & 
$((-4)^3)^5 = (-4)^{15}$ \\ \hline

\rowcolor{pinkcell!40}
Potencia de base racional & 
\centering $(a/b)^n = a^n / b^n$ & 
$(6/5)^2 = 6^2 / 5^2$ \\ \hline

\rowcolor{pinkcell!40}
Potencia de exponente negativo & 
\centering $a^{-n} = 1 / a^n$ & 
$8^{-3} = 1 / 8^3$ \\ \hline

\rowcolor{pinkcell!40}
Notación científica: operaciones & 
$a \cdot 10^{\pm n}$ siendo $1 \le a \le 9$.\newline
$+n$ para grandes números\newline
$-n$ para pequeños números & 
$320\,000\,000 = 3.2 \cdot 10^8$\newline
$0.000000009 = 9 \cdot 10^{-10}$ \\ \hline

\rowcolor{cyan!10}
Radicales: raíces de índice cualquiera & 
\multicolumn{2}{l|}{$\sqrt{49} = 7; \quad \sqrt[3]{-216} = -6; \quad \sqrt[4]{64} = 4; \quad \sqrt[4]{81} = 3; \quad \sqrt[5]{-32} = -2$} \\ \hline

\rowcolor{pinkcell!40}
Potencias de exponente racional & 
Una potencia con exponente racional puede expresarse en forma de raíz cuyo índice es el denominador del exponente y el radicando queda elevado al numerador del exponente: $x^{\frac{m}{n}} = \sqrt[n]{x^m}$ & 
$8^{2/5} = \sqrt[5]{8^2}$ \\ \hline

\rowcolor{cyan!10}
Extracción de factores de un radical & 
Si $m = n \cdot c + r$ entonces $\sqrt[n]{a^m} = a^c \cdot \sqrt[n]{a^r}$ & 
$\sqrt[3]{8^7} = 8^2 \cdot \sqrt[3]{8}$ \\ \hline

\rowcolor{pinkcell!40}
Operaciones con radicales & 
\centering $\sqrt[n]{x \cdot y \cdot z} = \sqrt[n]{x} \cdot \sqrt[n]{y} \cdot \sqrt[n]{z}$; \quad $\sqrt[n]{\dfrac{x}{y}} = \dfrac{\sqrt[n]{x}}{\sqrt[n]{y}}$ & 
$\sqrt[4]{5 \cdot 3 \cdot 2} = \sqrt[4]{5} \cdot \sqrt[4]{3} \cdot \sqrt[4]{2}$\newline
$\sqrt[3]{\dfrac{1}{27}} = \dfrac{\sqrt[3]{1}}{\sqrt[3]{27}} = \dfrac{1}{3}$ \\ \hline

\end{tabularx}
\end{center}

\end{document}